This calculation supplies the direct coefficient maps, positive metric,
traceless observable, finite polynomial, and integrated error for the
joint-Schur return. Its finite stopping constants are functions of the
specified packet and degree. Every matrix below is constructed from the
original sigma source and its full remainder map. Equation labels JRA1--JRA29
are local to this complete calculation.

\paragraph{Original objects and full coefficient maps.}
Retain $0<\delta<1/2$, $\gamma>2$, integer $m\geq1$, and
\begin{align}
 k&=4l+1,\quad l\geq1,\quad c=k/2,\quad e=1+k(m-1),\quad q=(k+1)^2e,\notag\\
 \lambda_{ab}&=c+(2a-k)\delta+i(2b-k)\gamma,\qquad
 \chi(S)=\prod_{a,b=0}^k(S-\lambda_{ab})^e,\qquad
 E=\C[S]/(\chi).\tag{JRA1}
\end{align}
The nodes are distinct: equality forces equality of their real and imaginary
parts, hence equality of $a,b$. None is real because $k$ is odd.
Put $\mathcal P_N=\{P\in\C[S]:\deg P\leq N\}$ with its given sigma form
\begin{equation}
 \langle P,Q\rangle_\sigma=\int_\R
 \overline{P(c+iy)}Q(c+iy)
 \frac{|\Gamma(1/4+iy/2)|^2}{2\pi}\,dy,
 \qquad H_N^\sigma=(\langle S^i,S^j\rangle_\sigma)_{i,j=0}^N.
 \tag{JRA2}
\end{equation}
These are the finite positive source forms of the incoming joint-Schur
construction. Positivity follows directly since their density is positive
and a nonzero polynomial cannot vanish on the entire line; finiteness is
the Gamma-moment property of that original source. We retain its mass
$c_\sigma=\sqrt{2\pi}$ in every $H_N^\sigma$.

The coefficient map $J_N:\mathcal P_N\to E$ is monic remainder modulo $\chi$,
in the basis $1,S,\ldots,S^{q-1}$. It is onto for $N\geq q-1$.
For $\alpha_j=c+2j+1/2$ and $f(S)=\prod_{j=0}^{l-1}(S-\alpha_j)$ define
$Z_M:\mathcal P_M\to\C^l$ by $(Z_MP)_j=P(\alpha_j)$.
For $M=N+l$ the joint map is
\begin{equation}
 \mathscr J_M=\binom{J_M}{Z_M}:\mathcal P_M\longrightarrow E\oplus\C^l,
 \quad
 \mathscr J_M(H_M^\sigma)^{-1}\mathscr J_M^*
 =\begin{pmatrix}K_M&C_M\\C_M^*&A_M\end{pmatrix}.
 \tag{JRA3}
\end{equation}
To prove surjectivity, $\chi$ and $f$ are coprime by JRA1 and the reality
of $\alpha_j$. Their Chinese remainder map is onto and has representatives
of degree at most $q+l-1\leq M$. Its two factors are $E$ and the evaluation
algebra at the distinct $\alpha_j$. Thus the block covariance in JRA3 is
positive definite: its quadratic form on $v\ne0$ is the strictly positive
$(H_M^\sigma)^{-1}$ form on $\mathscr J_M^*v\ne0$.
Consequently $K_M,A_M$ and both Schur complements
\begin{equation}
 G_M^\sigma=K_M^{-1},\qquad
 S_M=K_M-C_M A_M^{-1}C_M^*,\qquad
 L_M=A_M-C_M^*G_M^\sigma C_M
 \tag{JRA4}
\end{equation}
are positive definite. For example, minimizing the JRA3 form over the
first coordinate with second coordinate $z$ gives $z^*L_Mz>0$;
minimizing over the second proves the assertion for $S_M$.

Let $p_n$ be the monic orthogonal polynomials in this same sigma form,
with squared norms $\omega_n^\sigma>0$. Orthogonal decomposition of
$\mathcal P_{N+l}$ into $\mathcal P_N$ and its monic tail gives
\begin{align}
 U_N&=[Jp_{N+1},\ldots,Jp_{N+l}]:\C^l\to E,\quad
 \Omega_N=\diag(\omega_{N+1}^\sigma,\ldots,\omega_{N+l}^\sigma),\notag\\
 X_N&=\Omega_N^{-1}U_N^*G_N^\sigma U_N,\quad
 Y_N=A_{N+l}^{-1}C_{N+l}^*G_{N+l}^\sigma C_{N+l}.
 \tag{JRA5}
\end{align}
Thus $\Omega_NX_N=U_N^*G_N^\sigma U_N\geq0$ and
$A_{N+l}Y_N=C_{N+l}^*G_{N+l}^\sigma C_{N+l}\geq0$.
Moreover $A_{N+l}(I-Y_N)=L_{N+l}>0$.
The original metrics are $\Omega_N$ on the first coefficient space and
$A_{N+l}$ on the second. It follows, by conjugation with their positive
square roots, that $X_N$ has nonnegative real eigenvalues and $Y_N$ has
real eigenvalues in $[0,1)$. No cross covariance has been removed.

In raw derivative coordinates the remainder coordinate map is
\begin{equation}
\begin{gathered}
 (\mathsf E_\chi)_{(\lambda,r),n}
 =\begin{cases}n!\lambda^{n-r}/(n-r)!,&n\geq r,\\0,&n<r,\end{cases}
 \qquad
 \mathsf E_\chi M_f\mathsf E_\chi^{-1}
 =\bigoplus_\lambda F_\lambda,\\
 (F_\lambda)_{r,s}=
 \begin{cases}\binom rs f^{(r-s)}(\lambda),&r\geq s,\\0,&r<s.\end{cases}
\end{gathered}
 \tag{JRA6}
\end{equation}
All $0\leq r<e$ occur, and entries between different nodes are zero.
The equality is Leibniz's formula.
The inverse is Hermite interpolation, since vanishing of all these
derivatives is divisibility by $\chi$ and the domain has dimension $q$.
If an earlier file uses divided derivatives, its explicit map to JRA6 is
$D=\diag_{\lambda,r}(r!)$: raw jets equal $D$ times divided jets, and
the multiplication matrices obey $F_{\rm raw}=D F_{\rm div}D^{-1}$.
This does not change the original remainder map or drop a higher jet.

\paragraph{The determinant character representing the actual return.}
Write $V_N^\sigma=\det G_N^\sigma$ and let $G_N^0$ denote the tensor-Gamma
quotient form in the original remainder coordinates. Gamma recurrence on
$k=4l+1$ gives exactly
\begin{equation}
 dm_k^0=C_k|f(c+iy)|^2d\sigma(y),\qquad
 C_k=\frac{(\sqrt{2\pi})^k}{2^{1-k/2}\Gamma(k/2)}4^{-l}.
 \tag{JRA7}
\end{equation}
Indeed the holomorphic ratio is $(i/2)^l\prod_{j=0}^{l-1}(y-i(2j+1/2))$;
taking its real-line modulus gives JRA7, including the factor $4^{-l}$.
Multiplication by $f$ maps $\mathcal P_N$ isomorphically onto
$\ker Z_{N+l}$, sends its remainder $v$ to $M_fv$, and changes squared
norms by $C_k$. Minimizing the sigma norm under the joint constraint
$(M_fv,0)$, then inverting the positive block matrix JRA3, proves
\begin{equation}
 G_N^0=C_kM_f^*S_{N+l}^{-1}M_f,\qquad
 \log V_N^0=q\log C_k+2\log|\det M_f|
            +\log V_{N+l}^\sigma-\log\det(I-Y_N).
 \tag{JRA8}
\end{equation}
For the determinant equality, Schur elimination gives
$\det S_M\det A_M=\det K_M\det L_M$; divide by
$\det A_M\det K_M$ and use $I-Y_N=A_M^{-1}L_M$.
Here $M_f$ is invertible because $f,\chi$ are coprime.

The tail decomposition used in JRA5 gives
$K_{N+l}=K_N+U_N\Omega_N^{-1}U_N^*$.
Taking determinants and using $\det(I+AB)=\det(I+BA)$ gives
\begin{equation}
 \log V_N^\sigma-\log V_{N+l}^\sigma=\log\det(I+X_N).
 \tag{JRA9}
\end{equation}
The determinant identity follows, for instance, by computing the
determinant of $\left(\begin{smallmatrix}I&A\\-B&I\end{smallmatrix}\right)$
by either Schur complement.
Set
\begin{equation}
 (N_0,N_1,N_2,N_3)=(q-1,q,2q-1,2q),\quad
 (\varepsilon_0,\varepsilon_1,\varepsilon_2,\varepsilon_3)=(1,1,-1,-1).
 \tag{JRA10}
\end{equation}
The signed quotient-volume difference is therefore
\begin{equation}
 \mathcal R:=\sum_{r=0}^3\varepsilon_r
             (\log V_{N_r}^\sigma-\log V_{N_r}^0)
 =\sum_{r=0}^3\varepsilon_r
       \{\log\det(I+X_{N_r})+\log\det(I-Y_{N_r})\}.
 \tag{JRA11}
\end{equation}
The $q\log C_k+2\log|\det M_f|$ terms in JRA8 cancel because
$\sum_r\varepsilon_r=0$, after their four original appearances have
been retained. In the incoming notation, $\mathcal R=R_k^0-R_k^\sigma$.

Define the ordered direct sum and its metric by
\begin{align}
 \mathcal V&=\bigoplus_{r=0}^3
             (\mathcal V_{r,+}\oplus\mathcal V_{r,-}),\qquad
 \mathcal V_{r,\pm}=\C^l,\qquad n=8l,\notag\\
 \mathcal M&=\diag_{r=0}^3(\Omega_{N_r},A_{N_r+l}),\notag\\
 \mathcal B&=\diag_{r=0}^3(I+X_{N_r},I-Y_{N_r}),\qquad
 Q=\diag_{r=0}^3(\varepsilon_r I_l,\varepsilon_r I_l).
 \tag{JRA12}
\end{align}
Both $Q$ and $\mathcal B$ are $\mathcal M$-self-adjoint,
$\mathcal M\mathcal B>0$, and direct multiplication gives
\begin{equation}
 Q^2=I,\qquad \Tr Q=2l\sum_r\varepsilon_r=0,
 \qquad \Tr Q^2=n=8l.\tag{JRA13}
\end{equation}
The exact scalar map is the determinant character
\begin{equation}
 \Delta:\prod_{r,\pm}\operatorname{GL}(\mathcal V_{r,\pm})\to\C^\times,
 \qquad (T_{r,\pm})\longmapsto
       \prod_{r=0}^3\{\det T_{r,+}\det T_{r,-}\}^{\varepsilon_r}.
 \tag{JRA14}
\end{equation}
It is multiplicative because each determinant is. Equivalently it is the
action on $\bigotimes_r(\det\mathcal V_{r,+}\otimes
\det\mathcal V_{r,-})^{\varepsilon_r}$, with exponent $-1$ meaning the
dual line. On $\mathcal B$, its value is positive and
$\log\Delta(\mathcal B)=\mathcal R=\Tr(Q\log\mathcal B)$ by JRA11.
This is the typed map connecting the new eight coefficient summands to
the four original quotient determinants; it does not assert that $Q$
is the earlier four-projector contrast on the much larger source.

\paragraph{A complete simultaneous spectral frame.}
Construct $W$ as the direct sum of the positive square roots of the eight
blocks $\Omega_{N_r},A_{N_r+l}$ of $\mathcal M$.
Then $W^*W=\mathcal M$, $W^{-1}$ exists, and
$\widehat B=W\mathcal BW^{-1}$ is Hermitian positive.
Construct an orthonormal eigenbasis for $\widehat B$ separately in its
eight invariant blocks. For completeness, a Hermitian matrix has a
real eigenvalue by the fundamental theorem of algebra and
$\lambda v^*v=v^*Hv=\overline{v^*Hv}$; the orthogonal complement of an
eigenvector is invariant because $v^*Hw=(Hv)^*w$. Induction constructs
an orthonormal eigenbasis, and positivity makes all eigenvalues positive.
Apply the same construction to the positive matrix $\mathcal M$ to
construct its positive square root, taking the positive square root of
each eigenvalue. Thus both constructions and both inverse laws follow.
If $U$ is the resulting block unitary matrix, put $F=U^*W$.
Then, including every eigenvalue with multiplicity,
\begin{equation}
\begin{gathered}
 F^*F=\mathcal M,\qquad FF^{-1}=F^{-1}F=I,\\
 F\mathcal BF^{-1}=\diag(b_1,\ldots,b_n),\qquad b_i>0,\\
 FQF^{-1}=\diag(\epsilon_1,\ldots,\epsilon_n).
\end{gathered}
 \tag{JRA15}
\end{equation}
Here each $\varepsilon_r$ occurs $2l$ times. The last equality follows
because $Q$ is scalar on each of the eight blocks, so the constructed
block isometry commutes with its signs. Conjugation preserves products,
powers, inverses, and traces by direct multiplication and cyclicity.

\paragraph{Finite resolvent with explicit positive constants.}
Constants requiring only the eight original matrices are
\begin{align}
 \mu&=\min\bigl(1,\min_{0\leq r\leq3}\det(I-Y_{N_r})\bigr)>0,
 &\nu&=1+\sum_{r=0}^3\Tr X_{N_r}\geq1,\notag\\
 d&=(\mu+\nu)/2>0,
 &\theta&=(\nu-\mu)/(\nu+\mu)\in[0,1),\notag\\
 K_0&=\mu^{-2}\Tr(\mathcal B-I)^2
 =\mu^{-2}\sum_{r=0}^3\Tr(X_{N_r}^2+Y_{N_r}^2)\geq0.
 \tag{JRA16}
\end{align}
Every eigenvalue of $I-Y_N$ lies in $(0,1]$. Its determinant, the product
of all eigenvalues, is at most each one; hence $b_i\geq\mu$ on these
blocks. On the other blocks $b_i\geq1\geq\mu$ and
$b_i\leq1+\Tr X_{N_r}\leq\nu$. The minus blocks also have
$b_i\leq1\leq\nu$. This proves $\mu\leq b_i\leq\nu$ and all
claims in JRA16, including $\mu\leq1$ and $\mu\leq\nu$.

For $0\leq t\leq1$ set
\begin{align}
 B_t&=(1-t)I+t\mathcal B,&a_t&=1+t(d-1)>0,\notag\\
 X_t^{\rm res}&=B_t^{-1}(\mathcal B-I),&
 H_t&=I-B_t/a_t=t(dI-\mathcal B)/a_t,\notag\\
 T_t&=(\mathcal B-I)/a_t,&
 S_L(t)&=\sum_{j=0}^L H_t^jT_t\quad(L\geq0).
 \tag{JRA17}
\end{align}
All these endomorphisms are self-adjoint in the original $\mathcal M$
and commute with each other, being real rational functions of
$\mathcal B$ with positive denominators. In JRA15 their diagonal values are
\begin{equation}
 \lambda_i(t)=1+t(b_i-1),\quad
 \gamma_i(t)=\frac{b_i-1}{\lambda_i(t)},\quad
 h_i(t)=\frac{t(d-b_i)}{a_t}.
 \tag{JRA18}
\end{equation}
These equations explicitly supply the simultaneous diagonalization
required by the PR31 spectral-residual theorem. They follow by applying
$F$ to each factor of JRA17; the diagonal values have not been assigned
to an unrelated replacement matrix.
We have $\lambda_i(t)\geq\min(1,b_i)\geq\mu$ and
$|\gamma_i(t)|\leq|b_i-1|/\mu$.
Since $|d-b_i|\leq(\nu-\mu)/2$ and
$t/a_t\leq1/d$ (equivalently $td\leq1+t(d-1)$),
\begin{equation}
 |h_i(t)|\leq\theta<1,\qquad
 \sum_i\gamma_i(t)^2\leq K_0.
 \tag{JRA19}
\end{equation}

Direct multiplication gives $(I-H_t)X_t^{\rm res}=T_t$. The finite
telescoping identity
$\left(\sum_{j=0}^L H_t^j\right)(I-H_t)=I-H_t^{L+1}$, obtained by
cancelling every intermediate power, therefore proves
\begin{equation}
 R_L(t):=X_t^{\rm res}-S_L(t)=H_t^{L+1}X_t^{\rm res},\qquad
 E_L(t):=R_L(t)-\frac{\Tr R_L(t)}{n}I.
 \tag{JRA20}
\end{equation}
The constant term in $S_L$ and the exponent $L+1$ are both retained.

\paragraph{The actual centered pointwise error and its integral.}
Under $F$, the eigenvalues of $R_L(t)$ are
$\delta_i(t)=h_i(t)^{L+1}\gamma_i(t)$. Thus its trace is real and
\begin{align}
 V_L(t):=\Tr E_L(t)^2
 &=\sum_i\delta_i(t)^2-\frac{(\sum_i\delta_i(t))^2}{n}\notag\\
 &=\sum_i\left(\delta_i(t)-\frac1n\sum_j\delta_j(t)\right)^2\geq0,
 &V_L(t)&\leq\theta^{2L+2}K_0.
 \tag{JRA21}
\end{align}
Both equalities follow by expanding the square, and the bound follows
from JRA19 after dropping its nonnegative subtracted scalar square.
Scalar centering retains the signed pairing because $\Tr Q=0$:
$\Tr(R_LQ)=\Tr(E_LQ)$. For Hermitian matrices $A,C$ the entrywise
Cauchy--Schwarz inequality gives
$|\Tr(AC)|\leq(\Tr A^2\Tr C^2)^{1/2}$, since
$\Tr A^2=\sum_{i,j}|A_{ij}|^2$ and similarly for $C$.
Apply this to $FE_LF^{-1}$ and $FQF^{-1}$; trace invariance and JRA13 give
the full original-metric error
\begin{equation}
 \left|\Tr(QX_t^{\rm res})-\Tr(QS_L(t))\right|
 \leq\sqrt{8l\,V_L(t)}
 \leq\theta^{L+1}\sqrt{8l\,K_0}.
 \tag{JRA22}
\end{equation}
This reproduces the centered PR31 estimate on the specified joint-Schur
matrices, with its exact trace $8l$. It uses the generic signed trace
and spectral-residual results; it does not substitute $8l$ for the
different source-projector estimate $8q-6$.

For each $i$, elementary integration yields
$\log b_i=\int_0^1(b_i-1)/(1+t(b_i-1))\,dt$, also when $b_i=1$.
All denominators are at least $\mu>0$, so these finite sums are continuous
and can be integrated term by term. JRA11, JRA15 and JRA18 prove
\begin{equation}
\begin{gathered}
 \mathcal R=\int_0^1\Tr(QX_t^{\rm res})\,dt,\qquad
 \mathcal C_L=\int_0^1\Tr(QS_L(t))\,dt,\\
 |\mathcal R-\mathcal C_L|
 \leq\int_0^1\sqrt{8l\,V_L(t)}\,dt
 \leq\theta^{L+1}\sqrt{8l\,K_0}.
\end{gathered}
 \tag{JRA23}
\end{equation}
The inequality is integration of JRA22 on a unit interval, not a
claim that a pointwise bound automatically has been integrated.

The finite center itself has a closed expression in the original
matrix coefficients. Put $I_j(d)=\int_0^1t^j/a_t^{j+1}\,dt$. Then
\begin{align}
 \mathcal C_L
 &=\sum_{j=0}^L I_j(d)\Tr\bigl(Q(dI-\mathcal B)^j(\mathcal B-I)\bigr),\notag\\
 I_j(1)&=\frac1{j+1},\notag\\
 I_j(d)&=\frac{\displaystyle\log d+
       \sum_{a=1}^j(-1)^a\binom ja\frac{1-d^{-a}}a}
       {(d-1)^{j+1}}\qquad(d\ne1).
 \tag{JRA24}
\end{align}
For $d\ne1$, substitute $u=1+(d-1)t$; the integrand becomes
$(d-1)^{-j-1}(u-1)^ju^{-j-1}du$ on the oriented interval from $1$ to $d$.
Expand $(u-1)^j$ and integrate $u^{-1}$ and $u^{-a-1}$ to obtain JRA24,
including the case $0<d<1$ with its orientation. For $d=1$ the integrand
is $t^j$. Thus no unresolved matrix logarithm or path integral is needed
to compute the finite center. Every trace in it is the signed sum of
eight actual coefficient-matrix traces.

For $\eta>0$, put $A=\sqrt{8lK_0}$. If $A=0$, then JRA16 and positivity
of the eigenvalue squares force $\mathcal B=I$, so $\mathcal R=0$ and
$\mathcal C_0=0$. If $\theta=0$, JRA16 forces $\mu=\nu=1$, again
$\mathcal B=I$. Otherwise, with $A>0$ and $0<\theta<1$, the explicit
integer
\begin{equation}
 L_\eta=\max\left(0,
       \left\lceil\frac{\log(A/\eta)}{-\log\theta}\right\rceil\right)
 \quad\hbox{satisfies}\quad
 |\mathcal R-\mathcal C_{L_\eta}|\leq\eta.
 \tag{JRA25}
\end{equation}
Indeed, if $A\leq\eta$ every $L\geq0$ works. Otherwise
$L_\eta\geq\log(A/\eta)/(-\log\theta)$ and
$\theta^{L_\eta+1}A\leq\theta^{L_\eta}A\leq\eta$.
This proves an executable finite stopping rule for the given data.

\paragraph{Exact replacement and support maps.}
Section 6 of the incoming note gives its positive-tail interval
$[L_{\rm tail},U_{\rm tail}]$ for precisely JRA11. The present completed
calculation supplies the simultaneously valid intersection
\begin{equation}
 \mathcal R\in[L_{\rm tail},U_{\rm tail}]
 \cap[\mathcal C_L-\theta^{L+1}\sqrt{8lK_0},
       \mathcal C_L+\theta^{L+1}\sqrt{8lK_0}].
 \tag{JRA26}
\end{equation}
The intersection contains $\mathcal R$ because each interval does.
Unlike the positive-tail lower/upper sums, JRA24 forms its signed center
before the single centered error is bounded. This is an additional
finite certificate; no claim that its radius always beats the earlier
interval is needed or made.

The relation-valued section correction remains the same map.
Let $B_N$ be the columns of multiplication by $\chi$ into $\mathcal P_N$,
and $R_N^\sigma:E\to\mathcal P_N$ the sigma minimal section.
For the tensor-Gamma form $H_N^0$, minimization over its relation fibre
gives the actual section and exact difference
\begin{equation}
 R_N^0=R_N^\sigma-B_N(B_N^*H_N^0B_N)^{-1}B_N^*H_N^0R_N^\sigma,
 \quad
 R_N^0-R_N^\sigma\in\operatorname{Hom}(E,\chi\mathcal P_{N-q}).
 \tag{JRA27}
\end{equation}
To prove the formula, write every representative of $v$ as
$R_N^\sigma v+B_Nz$, expand its positive quadratic form, and complete
the square in $z$. $B_N$ has independent columns, so its Gram is
positive definite. At $N=q-1$ it is the zero-dimensional matrix and
the correction is zero. The weighted cyclic inclusion obeys
\begin{equation}
 \eta M_f=M_{f(\sum s_i)}\eta,\qquad
 \eta[P]=\upsilon_h^{\otimes k}P(\sum s_i)1.
 \tag{JRA28}
\end{equation}
This equality is multiplication of commuting elements in the original
tensor algebra; all full unit factors and nilpotents remain. Since
$f$ is coprime to $\chi$, its Bezout inverse on $E$ intertwines the same
way on the image of $\eta$. Thus the determinant calculation transports
the same source classes through explicitly invertible coefficient maps.

On any fixed retained face and outer label $\omega$, every coefficient
map $T$ above acts by $(\omega,x)\mapsto(\omega,Tx)$, and external absence
maps to external absence. The represented zero $(\omega,0)$ stays at
that label. Composition is exact since
$(\omega,T_2(T_1x))=(\omega,(T_2T_1)x)$.
These statements are maps of labelled coefficient sets; they do not
assign extra multiplicative structure to an arbitrary linear map.
Neither the direct sum nor the determinant character deletes a
present empty face or identifies it with $\tau$.

Finally, for the existing arithmetic correction
$\mathcal B_k^{\rm ar}=\mathcal B_k^0+\mathcal R+\delta_k^\sigma$,
the proven interval $\delta_k^\sigma\in[a_k,b_k]$ for the arithmetic
correction propagates literally:
\begin{equation}
 \mathcal B_k^{\rm ar}\in
 [\mathcal B_k^0+\mathcal C_L+a_k-\theta^{L+1}\sqrt{8lK_0},
  \mathcal B_k^0+\mathcal C_L+b_k+\theta^{L+1}\sqrt{8lK_0}].
 \tag{JRA29}
\end{equation}
This last implication is addition of the two proved intervals; the
arithmetic correction theorem itself remains in its original proof
files. Every new bound through JRA25 is proved here for the actual
joint-Schur matrices. A degree-dependent choice of $\mu,\nu,K_0$ has
not been replaced by a degree-independent constant. The next analytic
calculation is their behaviour, and that of the signed center, along
the original growing family.

\paragraph{Formal-source provenance.}
The generic results used and reproved above are
\begin{quote}\ttfamily
SplitZero.ResolventSeries.residual\_exact,\\
SplitZero.SignedTraceEnclosure.signed\_interval, and\\
SplitZero.SpectralResidual.signed\_error
\end{quote}
\noindent at PR31 head
\texttt{b4060b25c21f1a49a5e7730e9aff76d4c93c820d}.
The source files and full theorem-scope review accompany this TeX.
JRA15--JRA21 construct the required spectral equations rather than
assuming them as unnamed inputs. JRA23--JRA25 add the path integration
and its closed finite center. This application has a written proof;
it is not represented as an already compiled Lean declaration.
